\clearpage
\section*{ÖZET}
\label{ozettr}

Günümüzde REST API yapıları; web uygulamaları, mobil sistemler, mikro servis mimarileri ve bulut tabanlı platformlarda yaygın olarak kullanılmaktadır. API kullanımının artmasıyla birlikte güvenlik açıkları da modern yazılım sistemleri için önemli bir tehdit haline gelmiştir. Özellikle hatalı kimlik doğrulama, yetkilendirme eksiklikleri, IDOR/BOLA zafiyetleri, mass assignment problemleri ve güvenlik kontrollerinin yetersizliği; saldırganların sistem üzerinde yetki yükseltmesine, hassas verilere erişmesine ve kritik işlemleri gerçekleştirmesine neden olabilmektedir. Bu proje kapsamında geliştirilen ``API Fortress'' sistemi, REST API güvenlik açıklarının uygulamalı olarak öğrenilmesini amaçlayan eğitim tabanlı bir saldırı ve savunma simülasyon laboratuvarı olarak tasarlanmıştır. Sistem, öğrencilerin ve geliştiricilerin güvenlik açıklarını kontrollü bir ortamda deneyimleyebilmesini sağlamak amacıyla iki temel yapıdan oluşmaktadır: güvensiz API ortamı ve güvenli API ortamı. Güvensiz ortamda bilinçli olarak bırakılan zafiyetler üzerinden CTF benzeri hedeflere ulaşılması amaçlanırken, güvenli ortamda aynı açıkların nasıl kapatıldığı ve hangi savunma mekanizmalarıyla engellendiği gösterilmektedir. Proje kapsamında Mass Assignment ile yetki yükseltme, gizli header kullanılarak kimlik doğrulama atlama, parolasız girişe izin veren Broken Authentication mantık hatası, IDOR/BOLA ile başka kullanıcı verilerine erişim ve Broken Function Level Authorization ile yetkisiz kullanıcı silme gibi zafiyet senaryoları uygulanmıştır. Güvenli sürümde ise JWT tabanlı kimlik doğrulama, rol tabanlı yetkilendirme, rate limiting, regex tabanlı Web Application Firewall (WAF), güvenli parola hashleme ve giriş doğrulama kontrolleri geliştirilmiştir. Sistem geliştirme sürecinde Python programlama dili, Flask framework'ü, SQLAlchemy ORM yapısı ve Docker tabanlı çalışma ortamı kullanılmıştır. Gerçekleştirilen testler sonucunda API Fortress'in REST API güvenliği konusunda uygulamalı öğrenme sağlayan işlevsel bir laboratuvar ortamı sunduğu görülmüştür. Bu çalışma, güvenli yazılım geliştirme farkındalığını artırmayı ve API güvenliği alanında saldırı-savunma ilişkisini anlaşılır hale getirmeyi amaçlamaktadır.

\vspace{3mm}

\noindent \textbf{Anahtar Kelimeler:} REST API, API Güvenliği, OWASP API Security, CTF, Flask, JWT, WAF, IDOR, Mass Assignment, Docker, SQLAlchemy

\clearpage